\documentclass[11pt]{article}
\usepackage[margin=1in]{geometry}
\usepackage{amsmath,amssymb,amsthm}
\usepackage[hidelinks]{hyperref}
\usepackage{microtype}

\newcommand{\pv}{\operatorname{p.v.}}

\title{A Later Negative Answer to Wang--Zhang Problem 4\\
\large Smooth-phase Stein--Wainger principal values}
\author{FA\_BANACH\_001, agent lane 07}
\date{19 July 2026}

\begin{document}
\maketitle

\noindent\textbf{Status.}
\texttt{literature\_already\_answered} (exact negative answer).  This is a
literature-identification packet, not a claim of a new solution.

\section*{Original question}
Wang and Zhang ask in Section~4, Problem~4 (bottom of PDF page~14) of
\cite{WangZhang}:
\begin{quote}
Let $\lambda\in\mathbb R$ and $\psi\in C^\infty(\mathbb R)$.  Is
\[
 m_\psi(\lambda)=\pv\int_{-1}^{1}e^{i\lambda\psi(t)}\,\frac{dt}{t}
\]
bounded on $\mathbb R$?
\end{quote}

\section*{Decisive later answer}
Zhang and Zhu explicitly state that they resolve Wang--Zhang's problem
\cite[abstract and \S1.2]{ZhangZhu}.  Their Theorem~1.2 says that every real
$\psi\in C^\infty([-1,1])$ satisfies
\[
 m_\psi(\lambda)=o(\log\lambda),
\]
but this is sharp in the following strong sense: for every integer $k\ge2$
there exist a real smooth phase $\psi$ and a sequence $\lambda_n\to\infty$
such that
\[
 |m_\psi(\lambda_n)|\asymp
 \frac{\log\lambda_n}{\log^{(k)}\!\lambda_n}.
\]
In particular, $m_\psi$ is unbounded, so the answer to Problem~4 is
\emph{no}.  A smooth phase on the compact interval extends smoothly to
$\mathbb R$ without changing this integral.  The supporting paper also says
verbatim in its introduction that its results give a negative answer to
``Wang--Zhang [28, Problem 4 in Section 4].''  Thus the identification is
explicitly recognized by the supporting authors.

\section*{Scope and search evidence}
The later theorem settles exactly the scalar boundedness question selected
from the source paper; the source paper's other geometric restriction
problems are not addressed by this packet.  A bounded search on 19 July 2026
used arXiv ids 1606.09346 and 2603.23238, the exact displayed principal-value
formula, the source title and authors, and ``Stein--Wainger oscillatory
integral.''  It found the decisive 2026 arXiv preprint and an unanswered 2018
MathOverflow instance of the same bounded-interval question, but no earlier
separate paper explicitly resolving the Wang--Zhang formulation.

\section*{Human-review recommendation}
Confirm the source location and compare Theorem~1.2 and the explicit
identification in Section~1.2 of arXiv:2603.23238.  No mathematical inference
beyond smooth extension from $[-1,1]$ to $\mathbb R$ is needed.

\begin{thebibliography}{9}
\bibitem{WangZhang}
X. Wang and C. Zhang,
\emph{Sharp endpoint estimates for eigenfunctions restricted to submanifolds
of codimension 2}, Adv. Math. 386 (2021), 107835;
arXiv:1606.09346.

\bibitem{ZhangZhu}
C. Zhang and Z. Zhu,
\emph{Maximal growth of the Stein--Wainger oscillatory integral},
arXiv:2603.23238v1 (2026).
\end{thebibliography}

\end{document}
